\chapter{지울 수 없는 큰 공책}

우리 반 선생님은 칭찬할 일이 있을 때마다 \textbf{칭찬 스티커}를 한 장씩 주세요.
스티커를 열 장 모으면 급식 때 먼저 줄을 설 수 있어요.
그래서 스티커는 우리 반에서 돈처럼 쓰여요.
지우는 하준이에게 스티커 두 장을 주고 색종이를 받았어요.
서연이는 도윤이가 빌려 간 스티커 한 장을 돌려받았고요.

그러다 다툼이 생겼어요.

“내가 너한테 세 장 줬잖아!”

“아니야, 두 장이었어!”

누구 말이 맞을까요?
아무도 적어 두지 않았으니 알 수가 없어요.

\section{첫 번째 방법: 선생님의 공책}

선생님이 공책을 한 권 마련하셨어요.
스티커가 오갈 때마다 선생님께 말씀드리면 선생님이 공책에 적어요.

\begin{center}
\small
\begin{tabular}{@{}lll@{}}
\toprule
날짜 & 일어난 일 & \\
\midrule
3월 4일 & 지우 → 하준 & 스티커 2장 \\
3월 4일 & 도윤 → 서연 & 스티커 1장 \\
3월 5일 & 하준 → 민서 & 스티커 3장 \\
\bottomrule
\end{tabular}
\end{center}

이제 다툼이 생기면 공책을 펼쳐 보면 돼요.
그런데 걱정거리가 있어요.
선생님이 아파서 학교에 못 오시는 날에는 어떡하죠?
선생님이 실수로 잘못 적으시면요?
만약에, 아주 만약에 누군가 몰래 공책을 고친다면요?
공책이 한 권뿐이라면 그 공책을 믿는 수밖에 없어요.

\section{두 번째 방법: 모두의 공책}

그래서 우리 반은 다른 방법을 생각해 냈어요.
모든 학생이 똑같은 공책을 한 권씩 갖는 거예요.
스티커가 오가면 그 일을 큰 소리로 알리고, 모두가 자기 공책에 똑같이 적어요.

이제 누가 몰래 자기 공책을 고쳐도 소용없어요.
다른 스물다섯 권의 공책과 맞지 않으니 금방 들통나요.
선생님이 안 계셔도 공책은 스물여섯 권이나 남아 있어요.

한 가지 규칙을 더 정했어요.
공책은 한 쪽씩 채워 나가는데, 새 쪽 맨 위에는 앞쪽 내용을 요약한 \textbf{도장}을 찍어요.
도장은 앞쪽 글자가 하나만 달라져도 완전히 다른 모양이 되도록 만들어요.
그러면 옛날 쪽을 고치는 순간 다음 쪽의 도장과 맞지 않게 돼요.
쪽과 쪽이 사슬처럼 이어져 있는 셈이에요.

\begin{figure}[h]
\centering
\begin{tikzpicture}[page/.style={draw=inkblue, fill=white, line width=0.8pt, minimum width=24mm, minimum height=27mm, align=center, font=\scriptsize}]
  \foreach \i/\x/\t in {1/0/{지우→하준 2\\도윤→서연 1}, 2/3.1/{하준→민서 3\\서연→지우 1}, 3/6.2/{민서→도윤 2\\지우→서연 1}} {
    \node[page] (p\i) at (\x,0) {\\[2mm]\t};
    \node[font=\scriptsize\bfseries, color=inkblue] at ($(p\i.north)+(0,-0.3)$) {\i쪽};
  }
  \foreach \i/\j in {1/2, 2/3} {
    \draw[-{Stealth}, line width=1pt, color=sunline] (p\i.east) -- node[above, font=\tiny, color=black, align=center] {도장} (p\j.west);
  }
\end{tikzpicture}
\caption{쪽마다 앞쪽을 요약한 도장이 찍혀 있어서, 옛날 쪽을 고치면 사슬이 끊어져요.}
\end{figure}

축하해요.
여러분은 방금 \textbf{블록체인}의 가장 중요한 생각을 만들어 냈어요.
블록체인에서는 공책의 한 쪽을 \textbf{블록}이라고 부르고, 블록이 사슬(체인)처럼 이어져 있어서 블록체인이에요.
진짜 블록체인에서는 스물여섯 명의 학생 대신 전 세계 수많은 컴퓨터가 같은 공책을 나누어 갖고 있어요.

\section{이름표와 스티커 통}

진짜 블록체인에는 학생 이름이 없다고 1장에서 말했지요.
대신 모든 사람은 \textbf{지갑}을 가져요.
지갑은 스티커를 담아 두는 통이라고 생각하면 돼요.
통마다 이름표가 붙어 있는데, 이름표에는 이름 대신 \texttt{0x}로 시작하는 마흔두 글자짜리 \textbf{주소}가 적혀 있어요.
이 통에서 스티커를 꺼낼 수 있는 열쇠는 주인만 갖고 있어요.

새 통은 누구나 공짜로 얼마든지 만들 수 있어요.
한 사람이 통을 백 개 가질 수도 있지요.
이 사실은 5장에서 아주 중요해져요.

\section{여러 가지 스티커}

블록체인에는 스티커가 한 종류만 있지 않아요.
파란 스티커, 별 스티커, 하트 스티커처럼 종류가 아주 많아요.
이런 스티커를 \textbf{토큰}이라고 불러요.
토큰마다 제멋대로 규칙이 다르면 곤란하니까, 어른들은 “토큰은 이렇게 보내고, 이렇게 허락한다”는 공통 규칙을 정했어요.
그 규칙의 이름이 \textbf{ERC-20}이에요.
제 논문 제목에 들어 있는 바로 그 이름이에요.

\section{로봇도 통을 가져요}

블록체인의 통 가운데에는 사람이 아니라 \textbf{프로그램}이 주인인 통도 있어요.
교실로 말하면 \emph{문구점 로봇}이에요.
문구점 로봇은 정해진 규칙대로만 움직여요.
“파란 스티커 두 장을 받으면 연필 한 자루를 준다”처럼요.
스티커 종류를 바꿔 주는 \emph{교환 로봇}, 스티커를 빌려주는 \emph{대출 로봇}도 있어요.
블록체인에서는 이런 로봇을 \textbf{스마트 계약}이라고 부르고, 로봇이 하는 일을 사람들은 흔히 \textbf{앱}이라고 불러요.

이 책에서는 이제부터 블록체인 이야기를 교실 이야기로 바꿔서 할 거예요.
그래도 늘 기억해 주세요.
지갑은 이름표 붙은 스티커 통이고, 토큰은 스티커이고, 앱은 로봇이에요.

\begin{deeper}[진짜 용어]
\begin{itemize}
\item 비트코인은 2008년 사토시 나카모토라는 이름으로 발표된 글에서 시작되었어요. 그 사람이 누구인지는 아직도 아무도 몰라요.
\item 블록의 “도장”은 해시라는 계산으로 만들어요. 입력이 조금만 달라도 결과가 완전히 달라지는 계산이에요.
\item 주소는 \texttt{0x} 뒤에 16진수 40자리가 붙은 모양이에요.
\item ERC-20은 이더리움에서 토큰을 만드는 표준으로, 2015년 11월 파비안 포겔슈텔러와 비탈릭 부테린이 제안했어요.
\item 이 논문의 기록은 \textbf{아비트럼 원}이라는 블록체인에서 가져왔어요. 이더리움 위에서 더 빠르고 싸게 기록하도록 만든 블록체인이에요.
\end{itemize}
\end{deeper}

\begin{learned}
\begin{itemize}
\item 블록체인은 모두가 한 권씩 나누어 가진, 지울 수 없는 공책이에요.
\item 지갑은 이름표 붙은 스티커 통이고, 이름표에는 이름 대신 주소가 적혀 있어요.
\item 토큰은 스티커이고, ERC-20은 스티커를 다루는 공통 규칙이에요.
\item 사람이 아니라 프로그램(로봇)이 주인인 통도 있어요.
\end{itemize}
\end{learned}
